\begin{table}[t]
\centering
\caption{Long-horizon summary metrics at $2\,\tau_\lambda$ ($\approx 168$ steps) and valid prediction horizon (VPH). VPH uses the fixed full-test climatology convention of Fig.~1. HINE-L2 and MSR-HINE values are mean $\pm$ one standard deviation over 10 test trajectories; the VPH center is the crossing of the mean ACC curve and its uncertainty is the per-trajectory VPH spread. $\Delta$\% columns show relative improvement over HINE (positive\,=\,better). UNet-AR and FNO-AR diverge before $2\,\tau_\lambda$ (—).}
\label{tab:summary_metrics}
\resizebox{\columnwidth}{!}{%
\begin{tabular}{lrrrrrrrr}
\toprule
Model & VPH\,($\tau_\lambda$) & RMSE & ACC & SpecErr & $\Delta$RMSE\,\% & $\Delta$ACC\,\% & $\Delta$SpecErr\,\% \\
\midrule
UNet-AR & 0.155 & — & — & — & — & — & — \\
FNO-AR & 0.668 & 931.0254 & 0.0212 & — & — & -87.7 & — \\
HINE-L2 & 1.264 $\pm$ 0.172 & 6.1244 $\pm$ 0.6578 & 0.1717 $\pm$ 0.0546 & 0.2441 $\pm$ 0.0939 & +0.0 & -0.0 & +0.0 \\
MSR-HINE & 1.860 $\pm$ 0.139 & 4.8343 $\pm$ 0.5636 & 0.4387 $\pm$ 0.0562 & 0.1110 $\pm$ 0.0299 & +21.1 & +155.4 & +54.5 \\
\bottomrule
\end{tabular}}
\end{table}